\section*{Hoofdstuk \ref{ch:b2:metamorphic-facies} --- Metamorfe reacties en facies}

\begin{solution}{exo:b2:metamorphic-facies:1}
(a) Contactmetamorfose, een lokale thermische \omterm{def:b2:metamorphic-facies:metamorphism}{metamorfose} door de warmte van de intrusie.
(b) Dislocatiemetamorfose, in een breukzone: verkleining van de korrelgrootte, een myloniet.
(c) Oceaanbodemmetamorfose, door heet zeewater bij de rug, met onvolledige
herkristallisatie. (d) Begravingsmetamorfose, regionaal, zonder vervorming of
schistositeit.
\end{solution}

\begin{solution}{exo:b2:metamorphic-facies:2}
$\Delta_r S = 95.40 - 82.80 = \qty{12.60}{J.K^{-1}}$, $\Delta_r V = 4.986 - 4.415 =
\qty{0.571}{J.bar^{-1}}$; helling $12.60/0.571 = \qty{22.1}{bar/K} = \qty{2.21}{MPa/K}$.
Over \qty{100}{\celsius} stijgt de evenwichtsdruk met \qty{2.21}{kbar}.
\end{solution}

\begin{solution}{exo:b2:metamorphic-facies:3}
(a) Groenschist. (b) Blauwschist. (c) Amfiboliet: hoornblende met plagioklaas met meer
dan \qty{17}{\percent} anorthiet. (d) Eclogiet.
\end{solution}

\begin{solution}{exo:b2:metamorphic-facies:4}
(a) $\rho = M/V$: kyaniet $162.046/44.15 = \qty{3.670}{g.cm^{-3}}$, andalusiet
$162.046/51.52 = 3.145$, sillimaniet $162.046/49.86 = 3.250$. (b) Kyaniet: bij hogere
druk groeit de term $V\,(P - 1)$ van de gibbsenergie het minst voor het kleinste volume.
(c) Natuurlijke kristallen bevatten andere elementen (ijzer(III) bijvoorbeeld vervangt wat
aluminium in natuurlijk andalusiet en sillimaniet) en de metingen aan natuurlijke
exemplaren zijn minder precies dan de molaire volumes uit de tabellen.
\end{solution}

\begin{solution}{exo:b2:metamorphic-facies:5}
(a) $\Delta_r G = 3900 - 298.15\times8.59 = \qty{+1339}{J} > 0$: kyaniet. (b) Bij
\qty{1}{bar} geldt $\Delta_r H = T\Delta_r S$: $T = 3900/8.59 = \qty{454}{K} =
\qty{181}{\celsius}$. (c) $P = (8.59\times673.15 - 3900)/0.737 = 1882/0.737 =
\qty{2554}{bar} = \qty{2.55}{kbar}$, tegen 2.54 gemeten. (d) De reactie verloopt bij lage
temperatuur te traag: andalusiet blijft metastabiel bestaan, zoals de meeste
hogetemperatuurmineralen aan het oppervlak.
\end{solution}

\begin{solution}{exo:b2:metamorphic-facies:6}
(a) $\Delta_r H = 7700 + 1000\times(121.39 - 122.29) = \qty{6800}{J}$; $\Delta_r S = 296.44
- 285.02 = \qty{11.42}{J.K^{-1}}$; $\Delta_r G = 6800 - 1000\times11.42 = \qty{-4620}{J}$.
(b) Met de waarden bij \qty{298.15}{K}: $7700 - 1000\times12.60 = \qty{-4900}{J}$.
(c) Sillimaniet. De eenvoudige methode zit er \qty{280}{J} naast, ongeveer \qty{6}{\percent}; gedeeld
door $\Delta_r V = 0.571$ is dat ongeveer \qty{490}{bar} op de lijn, veel minder dan het
effect van de onzekerheid van de enthalpieën.
\end{solution}

\begin{solution}{exo:b2:metamorphic-facies:7}
(a) $\Delta_r G = -1127.4 + 1128.5 = \qty{1.1}{kJ}$ en $\Delta_r V = (34.15 -
36.93)/10 = \qty{-0.278}{J.bar^{-1}}$: $P = 1100/0.278 = \qty{3957}{bar} \approx
\qty{4.0}{kbar}$. (b) $\Delta_r S = 88.0 - 91.7 = -3.7$; helling $-3.7/-0.278 =
\qty{13.3}{bar/K}$. (c) Met $\Delta_r H = 0$ geldt $P = T\Delta_r S/\Delta_r V = 13.31\,T$:
bij \qty{573.15}{K} \qty{7629}{bar} $= \qty{7.63}{kbar}$. (d) Schelparagoniet is
metastabiel aan het oppervlak: organismen maken het, en het blijft bestaan omdat het maar
langzaam in calciet overgaat.
\end{solution}

\begin{solution}{exo:b2:metamorphic-facies:8}
(a) $-32.4/-1.698 = \qty{19.1}{bar/K}$. (b) $P = (T\Delta_r S - \Delta_r H)/\Delta_r V =
(-32.4\times773.15 + 5000)/(-1.698) = (-25050 + 5000)/(-1.698) = \qty{11808}{bar}
\approx \qty{11.8}{kbar}$. (c) $\sqrt{2.6^2 + 3.6^2 + 1.0^2} = \qty{4.55}{kJ}$ op
$\Delta_r H$, dus $4550/1.698 = \qty{2.7}{kbar}$ op de druk. (d) Bij de drukken van
de eclogietfacies is de albietcomponent van plagioklaas instabiel ten opzichte van
jadeïet en kwarts; het natrium gaat in het pyroxeen (omfaciet), en plagioklaas
verdwijnt.
\end{solution}

\begin{solution}{exo:b2:metamorphic-facies:9}
(a) Van 1 tot \qty{2}{kbar}: $1000/34 = \qty{29.4}{bar/K}$; van 2 tot 3: $1000/25 =
\qty{40.0}{bar/K}$. (b) Met vaste $\Delta_r S$ is $\Delta_r V$ omgekeerd evenredig met
de helling: het kromp met een factor $40.0/29.4 = 1.36$. (c) Het vrijgekomen water: zijn molaire
volume daalt naarmate de druk stijgt, terwijl de volumes van de vaste stoffen nauwelijks veranderen.
\end{solution}

\begin{solution}{exo:b2:metamorphic-facies:10}
(a) Bij \qty{873.15}{K}: kyaniet $=$ andalusiet $(8.59\times873.15 - 3900)/0.737 =
\qty{4885}{bar}$; kyaniet $=$ sillimaniet $(12.60\times873.15 - 7700)/0.571 =
\qty{5782}{bar}$; andalusiet $=$ sillimaniet $(4.01\times873.15 - 3800)/(-0.166) =
\qty{1799}{bar}$. (b) Bij \qty{4885}{bar}, onder de lijn kyaniet $=$ sillimaniet (5782),
is sillimaniet stabieler dan kyaniet, en boven de lijn andalusiet $=$ sillimaniet
(1799) stabieler dan andalusiet: het punt ligt in het sillimanietveld, en het
evenwicht kyaniet $=$ andalusiet is daar metastabiel. (c) Bij \qty{4}{kbar}: tussen 1799
en \qty{5782}{bar}, sillimaniet. Bij \qty{7}{kbar}: boven \qty{5782}{bar} wint kyaniet van
sillimaniet, en boven \qty{4885}{bar} van andalusiet: kyaniet.
\end{solution}

\begin{solution}{exo:b2:metamorphic-facies:11}
(a) $z = P/(\rho g) = 3.75\times10^8/(2700\times9.81) = \qty{14.2}{km}$, onzeker tot op
$0.25\times10^8/26487 = \qty{0.94}{km}$. (b) $(504 - 10)/14.16 =
\qty{35}{\celsius.km^{-1}}$. (c) De drie mineralen bestaan alleen in het \omterm{def:b2:metamorphic-facies:invariant}{invariante punt}
in evenwicht naast elkaar: het pad van het gesteente moet erdoorheen gaan, en de reacties moeten
daar stoppen. De meeste gesteenten gaan er langs en houden één of twee polymorfen over.
\end{solution}

\begin{solution}{exo:b2:metamorphic-facies:12}
(a) X: $T = 10 + 10\times30 = \qty{310}{\celsius}$ (\qty{583.15}{K}); $P = 3000\times
9.81\times30000 = \qty{8.83e8}{Pa} = \qty{8.83}{kbar}$. Y: $T = 10 + 30\times20 =
\qty{610}{\celsius}$ (\qty{883.15}{K}); $P = 2800\times9.81\times20000 = \qty{5.49}{kbar}$.
(b) X: calciet $=$ aragoniet bij $13.31\times583.15 = \qty{7.76}{kbar} < 8.83$, aragoniet;
albiet $=$ jadeïet + kwarts bij $(32.4\times583.15 - 5000)/1.698 = \qty{8.18}{kbar} <
8.83$, jadeïet + kwarts; kyaniet $=$ andalusiet bij $(8.59\times583.15 - 3900)/0.737 =
\qty{1.51}{kbar}$, onder de temperatuur van het tripelpunt: kyaniet. Y: calciet $=$
aragoniet bij \qty{11.75}{kbar} $> 5.49$, calciet; albiet $=$ jadeïet + kwarts bij
\qty{13.9}{kbar}, albiet; boven het tripelpunt kyaniet $=$ sillimaniet bij
\qty{6.00}{kbar} $> 5.49$ en andalusiet $=$ sillimaniet bij \qty{1.56}{kbar} $< 5.49$:
sillimaniet. (c) X, bij \qty{10}{\celsius\per\km}, is het koudere pad: jadeïet en
aragoniet bij slechts \qty{310}{\celsius} plaatsen het in de blauwschistfacies. Y, drie keer
warmer, eindigt met sillimaniet bij \qty{610}{\celsius} en matige druk, in de
amfibolietfacies.
\end{solution}

\begin{solution}{pb:b2:metamorphic-facies:1}
\noindent\textbf{Deel I.}

\textbf{1.} $\Delta_r H = -2589.9 + 2593.8 = \qty{3.9}{kJ}$; $\Delta_r S = 91.39 - 82.80 =
\qty{8.59}{J.K^{-1}}$; $\Delta_r V = 5.152 - 4.415 = \qty{0.737}{J.bar^{-1}}$.

\textbf{2.} $\qty{7.7}{kJ}$; $\qty{12.60}{J.K^{-1}}$; $\qty{0.571}{J.bar^{-1}}$.

\textbf{3.} $\qty{3.8}{kJ}$; $\qty{4.01}{J.K^{-1}}$; $\qty{-0.166}{J.bar^{-1}}$; elk is
de tweede min de eerste, omdat andalusiet wegvalt.

\textbf{4.} Sillimaniet heeft de grootste entropie (95.40): het ligt aan de
hogetemperatuurkant van zijn beide lijnen. Kyaniet heeft het kleinste volume (4.415): het ligt aan de
hogedrukkant van zijn beide lijnen. Andalusiet, met het grootste volume, beheerst de
lagedrukkant.

\textbf{5.} $8.59/0.737 = \qty{11.66}{bar/K}$; $12.60/0.571 = \qty{22.07}{bar/K}$;
$4.01/(-0.166) = \qty{-24.16}{bar/K}$.

\textbf{6.} Sillimaniet heeft meer entropie maar minder volume dan andalusiet: $\Delta_r S > 0$
en $\Delta_r V < 0$. Een hogere temperatuur bevoordeelt sillimaniet, een hogere druk
ook, dus de lijn loopt bij hogere temperatuur naar lagere druk.

\noindent\textbf{Deel II.}

\textbf{7.} $3900 - 298.15\times8.59 = \qty{+1339}{J}$; $7700 - 298.15\times12.60 =
\qty{+3943}{J}$. Beide positief: kyaniet is stabiel aan het oppervlak.

\textbf{8.} $T = \Delta_r H/\Delta_r S$: $3900/8.59 = \qty{454.0}{K} = \qty{181}{\celsius}$;
$7700/12.60 = \qty{611.1}{K} = \qty{338}{\celsius}$; $3800/4.01 = \qty{947.6}{K} =
\qty{674}{\celsius}$.

\textbf{9.} Bij verhitting bij \qty{1}{bar} maakt kyaniet plaats voor andalusiet bij
\qty{181}{\celsius}; bij \qty{338}{\celsius} is kyaniet niet meer stabiel, dus kyaniet $=$
sillimaniet is daar metastabiel; andalusiet maakt plaats voor sillimaniet bij
\qty{674}{\celsius}. Het eerste en het derde zijn stabiele evenwichten.

\textbf{10.} $(8.59\times673.15 - 3900)/0.737 = \qty{2554}{bar} = \qty{2.55}{kbar}$,
\qty{0.01}{kbar} van het experiment.

\textbf{11.} $(12.60\times973.15 - 7700)/0.571 = 4562/0.571 = \qty{7989}{bar} =
\qty{7.99}{kbar}$, \qty{0.13}{kbar} boven het experiment.

\textbf{12.} Kyaniet $=$ andalusiet stijgt vanaf \qty{181}{\celsius} bij \qty{1}{bar} met
helling \qty{11.7}{bar/K}; andalusiet $=$ sillimaniet daalt van ongeveer \qty{11.5}{kbar} bij
\qty{200}{\celsius} tot \qty{1}{bar} bij \qty{674}{\celsius}; kyaniet $=$ sillimaniet, de
steilste van de twee stijgende lijnen, snijdt beide bij ongeveer \qty{514}{\celsius} en \qty{3.9}{kbar}, waarbij elke lijn maar aan
één kant van het snijpunt stabiel is (het diagram van het hoofdstuk).

\noindent\textbf{Deel III.}

\textbf{13.} $a_1 = 8.59/0.737 = \qty{11.655}{bar/K}$, $b_1 = 3900/0.737 =
\qty{5291.7}{bar}$; $a_2 = 12.60/0.571 = \qty{22.067}{bar/K}$, $b_2 = 7700/0.571 =
\qty{13485.1}{bar}$.

\textbf{14.} $T = (b_2 - b_1)/(a_2 - a_1) = 8193.4/10.412 = \qty{786.92}{K} =
\qty{513.8}{\celsius}$.

\textbf{15.} $P = 11.655\times786.92 - 5291.7 = \qty{3880}{bar} = \qty{3.88}{kbar}$.

\textbf{16.} $(4.01\times786.92 - 3800)/(-0.166) = (3155.5 - 3800)/(-0.166) =
\qty{3882}{bar}$: dezelfde druk binnen de afronding van de gegevens, zoals de propositie van het hoofdstuk vereist.

\textbf{17.} \qty{514}{\celsius} ligt \qty{10}{\celsius} van \qty{504}{\celsius}, binnen
de onzekerheid van \qty{20}{\celsius}; \qty{3.88}{kbar} ligt \qty{0.13}{kbar} van 3.75,
binnen 0.25. De overeenstemming is goed.

\textbf{18.} Nee: \qty{536}{\celsius} ligt \qty{32}{\celsius} van het experiment, buiten de
onzekerheid ervan. Het verwaarlozen van de warmtecapaciteiten was niet de belangrijkste fout: de enthalpieën zijn dat.

\noindent\textbf{Deel IV.}

\textbf{19.} $\Delta_r H$ van kyaniet $\to$ andalusiet stijgt met \qty{1000}{J}; bij vaste
$T$ daalt $P = (T\Delta_r S - \Delta_r H)/\Delta_r V$ met $1000/0.737 = \qty{1357}{bar}$.

\textbf{20.} $\Delta T = -1357/(a_2 - a_1) = -1357/10.41 = \qty{-130}{K}$: het tripelpunt
verschuift naar ongeveer \qty{384}{\celsius}.

\textbf{21.} De lijn kyaniet $=$ sillimaniet daalt met $1000/0.571 = \qty{1751}{bar}$, en
het tripelpunt verschuift met $+1751/10.41 = \qty{+168}{K}$.

\textbf{22.} Elke enthalpie afzonderlijk, gevarieerd met $\pm\qty{2.0}{kJ}$, verschuift het tripelpunt met
$\pm\qty{260}{K}$ (andalusiet) of $\pm\qty{336}{K}$ (sillimaniet): de calorimetrie alleen laat
tripelpunten over enkele honderden graden toe.

\textbf{23.} De getabelleerde enthalpieën zijn niet zuiver calorimetrisch: binnen hun
onzekerheden zijn ze zo aangepast dat de berekende lijnen bij de omgekeerde experimenten passen.
De experimenten leggen het diagram vast; de tabellen dragen het, consistent, over naar andere reacties.

\textbf{24.} Het tripelpunt, berekend uit de tabellen met constante eigenschappen, ligt op
\textbf{\qty[text-series-to-math=true]{514}{\celsius}}.
\end{solution}
